\section{Motivation}
\label{sec:motivation}

Finding \emph{where a radio transmission is coming from} is a classical problem
with modern urgency: locating a lost drone by its video downlink, finding an
interference source that is degrading a shared band, tracking wildlife tags,
or navigating toward a beacon when GPS is unavailable. Commercial
direction-finding gear solves this with large, calibrated, many-element
antenna arrays and correspondingly large price tags.

SPF starts from the opposite end of the cost spectrum and asks:

\begin{quote}
\emph{How much localization capability can be extracted from the cheapest
possible sensor --- a single pair of antennas attached to a \$100--\$200
software-defined radio --- if we are willing to move the sensor, collect a lot
of real data, and learn the sensor's imperfections instead of engineering them
away?}
\end{quote}

Three observations make this question worth asking.

\paragraph{1. Two antennas are the minimal interferometer.}
A two-element array measures exactly one number per snapshot --- the phase
difference between its channels --- and that number is a (nonlinear,
ambiguous, noisy) function of the emitter bearing
(Section~\ref{sec:math}). Everything the platform will ever know about
direction must be squeezed out of this one degree of freedom, plus motion,
plus time. It is the sharpest possible test bed for ``sensing with learned
models'': the physics is simple enough to write down exactly, yet the
real-world measurement deviates from that physics in ways that are hard to
calibrate by hand (mount rotations, effective antenna spacing, gain-dependent
phase inversions, automatic gain control, multipath).

\paragraph{2. Motion converts ambiguity into information.}
A stationary two-element array cannot distinguish front from back, and for
wide antenna spacing cannot even pin down a unique bearing
(Section~\ref{sec:ambiguity}). But a \emph{moving} receiver-- or a pair of
receivers --- sweeps its ambiguity structure through space; a filter that
accumulates measurements over a trajectory can collapse the ambiguity. This
trades hardware complexity for algorithmic complexity, which is exactly the
trade a learned system should win.

\paragraph{3. Real data beats simulated calibration.}
Early experiments made it clear that the dominant error sources are not
thermal noise but \emph{systematics}: antennas mounted a few millimetres from
their nominal positions, arrays labeled with the wrong spacing, compass biases
that drift between calibration sessions, receiver gain control interacting
with the RF front end. These are properties of each physical rig on each
particular day. SPF therefore treats large-scale, ground-truth-annotated data
collection --- and the auditing of that data
(Section~\ref{sec:audit-approach}) --- as a first-class part of the project,
on equal footing with model architecture.
